\chapter{The Maximum Principle}

\section*{3.1 \quad Statement of the maximum principle}

\begin{theorem}
\label{topping-maxpr-weak-maximum-principle}
\source{topping:page-0045,topping:page-0046}
\dcref{ch3:1.1}
\group{maximum-principle}
\lean{Topping.weak_maximum_principle}
\leanok
Let $g(t)$ and $X(t)$ be smooth time-dependent metrics and vector fields on a
closed manifold $\M$, for $t\in[0,T]$ with $0<T<\infty$. Let
$F:\RR\times[0,T]\to\RR$ be smooth, and suppose
$u\in C^\infty(\M\times[0,T],\RR)$ satisfies
\[
\frac{\partial u}{\partial t}
\leq \Delta_{g(t)}u+\langle X(t),\nabla u\rangle+F(u,t).
\]
If $\phi:[0,T]\to\RR$ solves
\[
\begin{cases}
\dfrac{d\phi}{dt}=F(\phi(t),t),\\
\phi(0)=\alpha\in\RR,
\end{cases}
\]
and $u(\cdot,0)\leq\alpha$, then $u(\cdot,t)\leq\phi(t)$ for every
$t\in[0,T]$.
\end{theorem}

\begin{proof}
\sketch For $\e>0$, perturb the comparison solution by solving
\[
\begin{cases}
\dfrac{d\phi_\e}{dt}=F(\phi_\e(t),t)+\e,\\
\phi_\e(0)=\alpha+\e.
\end{cases}
\]
For sufficiently small $\e$, this solution exists on $[0,T]$ and converges
uniformly to $\phi$. If $u<\phi_\e$ first fails at $(x,t_0)$, then $x$ is a
spatial maximum at time $t_0$, so $\nabla u(x,t_0)=0$, $\Delta u(x,t_0)\leq0$,
and $\partial_tu(x,t_0)-\phi_\e'(t_0)\geq0$. The differential inequality and
the perturbed ODE then imply $0\geq\e>0$, a contradiction. Letting
$\e\downarrow0$ proves the claim.
\end{proof}

\begin{corollary}
\label{topping-maxpr-weak-minimum-principle}
\source{topping:page-0045}
\dcref{ch3:1.2}
\group{maximum-principle}
\lean{Topping.weak_minimum_principle}
\leanok
Theorem~\ref{topping-maxpr-weak-maximum-principle} remains valid after reversing all inequalities.
\end{corollary}

\begin{remark}
\label{topping-maxpr-strong-maximum-principle}
\source{topping:page-0046}
\dcref{ch3:1.3}
\group{maximum-principle}
\lean{Topping.eq_of_metricLaplacianAt_nonneg_of_isMaxOn,
Topping.lt_or_eqOn_of_metricLaplacianAt_nonneg}
The scalar strong maximum principle strengthens the conclusion to
$u(\cdot,t)<\phi(t)$ for $t\in(0,T]$, unless
$u(x,t)=\phi(t)$ for every $x\in\M$ and $t\in[0,T]$.
\end{remark}

\section*{3.2\ \ \ Basic control on the evolution of curvature}

\begin{theorem}
\label{topping-maxpr-scalar-lower-bound}
\source{topping:page-0047}
\uses{topping-maxpr-weak-minimum-principle}
\dcref{ch3:2.1}
\group{maximum-principle}
\lean{Topping.scalarLowerBarrier_le_of_scalarCurvatureEvolution}

Let $g(t)$ be a Ricci flow on a closed manifold $\M$ for $t\in[0,T]$. If
$R\geq\alpha\in\RR$ at $t=0$, then
\[
R\geq\frac{\alpha}{1-\left(\frac{2\alpha}{n}\right)t}
\]
for all $t\in[0,T]$.
\end{theorem}

\begin{proof}
\sketch Apply the weak minimum principle to the scalar curvature evolution
equation (2.5.6), with $u=R$, $X=0$, and $F(r,t)=\frac{2}{n}r^2$. Its
comparison ODE has solution
\[
\phi(t)=\frac{\alpha}{1-\left(\frac{2\alpha}{n}\right)t}.
\]
\end{proof}

\begin{corollary}
\label{topping-maxpr-scalar-preserved}
\source{topping:page-0047}
\dcref{ch3:2.2}
\group{maximum-principle}
\lean{Topping.scalarCurvature_ge_of_initial_ge}
For a Ricci flow $g(t)$ on a closed manifold $\M$, if $R\geq\alpha\in\RR$ at
$t=0$, then $R\geq\alpha$ throughout the time interval of the flow.
\end{corollary}

\begin{corollary}
\label{topping-maxpr-positive-scalar-preserved}
\source{topping:page-0047}
\dcref{ch3:2.3}
\group{maximum-principle}
\lean{Topping.scalarCurvature_pos_of_initial_pos,
Topping.scalarCurvature_nonneg_of_initial_nonneg}
Under Ricci flow on a closed manifold, positive and weakly positive scalar
curvature are preserved.
\end{corollary}

\begin{corollary}
\label{topping-maxpr-finite-time-blowup}
\source{topping:page-0047}
\dcref{ch3:2.4}
\group{maximum-principle}
\lean{Topping.le_of_scalarCurvature_initial_pos_Ico}
Let $g(t)$ be a Ricci flow on a closed manifold for $t\in[0,T)$. If
$R\geq\alpha>0$ at $t=0$, then $T\leq\frac{n}{2\alpha}$.
\end{corollary}

\begin{corollary}
\label{topping-maxpr-negative-scalar-lower-bound}
\source{topping:page-0047}
\dcref{ch3:2.5}
\group{maximum-principle}
\lean{Topping.neg_div_le_scalarCurvature}
For a Ricci flow on a closed manifold, defined for $t\in(0,T]$,
\[
R\geq-\frac{n}{2t}
\]
for every $t\in(0,T]$.
\end{corollary}

\begin{corollary}
\label{topping-maxpr-volume-decreases}
\source{topping:page-0047}
\dcref{ch3:2.6}
\group{maximum-principle}
\lean{Topping.riemannianVolume_antitoneOn_of_isRicciFlowOn_Icc}
\leanok
If a Ricci flow $g(t)$ on a closed manifold is defined for $t\in[0,T]$ and has
$R\geq0$ at $t=0$, then its volume $V(t)$ is weakly decreasing on $[0,T]$.
\end{corollary}

\begin{corollary}
\label{topping-maxpr-volume-ratio}
\source{topping:page-0047,topping:page-0048}
\uses{topping-maxpr-scalar-lower-bound}
\dcref{ch3:2.7}
\group{maximum-principle}
\lean{Topping.volumeRatio_antitoneOn,
Topping.volume_le_of_scalarCurvature_initial_ge}

Let $g(t)$ be a Ricci flow on a closed manifold for $t\in[0,T]$, and let
$\alpha:=\inf R<0$ at $t=0$. Then
\[
\frac{V(t)}{\left(1+\frac{2(-\alpha)}{n}t\right)^{\frac n2}}
\]
is weakly decreasing, and consequently
\[
V(t)\leq V(0)\left(1+\frac{2(-\alpha)}{n}t\right)^{\frac n2}.
\]
\end{corollary}

\begin{proof}
\sketch Differentiate the logarithm of the ratio. The result is
\[
\frac{\dd}{\dd t}\ln\!\left[
\frac{V(t)}{\left(1+\frac{2(-\alpha)}{n}t\right)^{\frac n2}}
\right]
=-\frac{1}{V}\int R\,\dd V
 +\frac{\alpha}{1+\frac{2(-\alpha)}{n}t}.
\]
This is at most
\[
-\inf_{\M}R(t)+\frac{\alpha}{1-\frac{2\alpha}{n}t}\leq0,
\]
by Theorem~\ref{topping-maxpr-scalar-lower-bound}.
\end{proof}

\begin{proposition}
\label{topping-maxpr-curvature-norm-evolution}
\source{topping:page-0048,topping:page-0049}
\dcref{ch3:2.10}
\group{maximum-principle}
\lean{Topping.HasCurvatureNormEvolutionInequalityOn,
Topping.exists_curvatureNormEvolution_const_of_ingredients,
Topping.exists_uniform_curvatureNormEvolution_const,
Topping.exists_uniform_curvatureNormEvolution_const_of_morganTian_isRicciFlowOn}
Under Ricci flow,
\[
\frac{\partial}{\partial t}|\Rm|^2
\leq\Delta|\Rm|^2-2|\nabla\Rm|^2+C|\Rm|^3,
\]
where $C=C(n)$ and $\Delta$ is the Laplace--Beltrami operator.
\end{proposition}

\begin{proof}
\sketch Differentiating the metric contraction defining $|\Rm|^2$ and using
the curvature evolution equation gives
\[
\frac{\partial}{\partial t}|\Rm|^2
\leq2\langle\Rm,\Delta\Rm\rangle+C|\Rm|^3.
\]
The identity
\[
\Delta|\Rm|^2=2|\nabla\Rm|^2+2\langle\Rm,\Delta\Rm\rangle
\]
then yields the displayed inequality after weakening the resulting estimate.
\end{proof}

\begin{theorem}
\label{topping-maxpr-curvature-norm-bound}
\source{topping:page-0049}
\dcref{ch3:2.11}
\group{maximum-principle}
\lean{Topping.riemannNormAt_le,
Topping.exists_uniform_riemannNormAt_le_of_ingredients,
Topping.exists_uniform_riemannNormAt_le_of_morganTian_isRicciFlowOn}
Let $g(t)$ be a Ricci flow on a closed manifold $\M$ for $t\in[0,T]$, and
suppose $|\Rm|\leq M$ at $t=0$. Then, for $t\in(0,T]$,
\[
|\Rm|\leq\frac{M}{1-\frac12CMt}.
\]
\end{theorem}

\begin{proof}
\sketch Apply the weak maximum principle to the weakened inequality from
Proposition~\ref{topping-maxpr-curvature-norm-evolution} with $u=|\Rm|^2$, $X=0$, $F(r,t)=Cr^{3/2}$, and initial
value $\alpha=M^2$. The comparison solution is
\[
\phi(t)=\frac{1}{\bigl(M^{-1}-\frac12Ct\bigr)^2},
\]
so $|\Rm|^2\leq\phi(t)$ and the stated bound follows.
\end{proof}

\section*{3.3\ \ \ Global curvature derivative estimates}

\begin{theorem}
\label{topping-maxpr-first-derivative-estimate}
\source{topping:page-0050,topping:page-0051}
\dcref{ch3:3.1}
\group{maximum-principle}
\lean{Topping.mul_le_of_gradient_parabolic_inequalities,
Topping.gradRiemannNorm_le_of_evolution_inequalities,
Topping.shiCoeff,
Topping.shi_cancel_sum_le,
Topping.shi_parabolic_inequality}
Let $M>0$ and let $g(t)$ be a Ricci flow on a closed manifold $\M^n$ for
$t\in[0,\frac1M]$. For every $k\in\NN$ there is $C=C(n,k)$ such that, if
$|\Rm|\leq M$ on $\M\times[0,\frac1M]$, then
\[
|\nabla^k\Rm|\leq\frac{CM}{t^{k/2}}
\]
for all $t\in[0,\frac1M]$.
\end{theorem}

\begin{proof}
\sketch For $k=1$, the curvature evolution and commutation formulae give
\[
\frac{\partial}{\partial t}\Rm=\Delta\Rm+\Rm*\Rm,
\qquad
\nabla(\Delta\Rm)=\Delta(\nabla\Rm)+\Rm*\nabla\Rm,
\]
and
\[
\nabla\frac{\partial}{\partial t}\Rm
=\frac{\partial}{\partial t}\nabla\Rm+\Rm*\nabla\Rm.
\]
Consequently,
\[
\frac{\partial}{\partial t}\nabla\Rm
=\Delta(\nabla\Rm)+\Rm*\nabla\Rm.
\]
Consequently,
\[
\frac{\partial}{\partial t}|\nabla\Rm|^2
\leq\Delta|\nabla\Rm|^2-2|\nabla^2\Rm|^2
 +C|\Rm|\,|\nabla\Rm|^2.
\]
For $u=t|\nabla\Rm|^2+\alpha|\Rm|^2$, the preceding inequality and
Proposition~\ref{topping-maxpr-curvature-norm-evolution} imply, for sufficiently large $\alpha$,
\[
\frac{\partial u}{\partial t}\leq\Delta u+C(n)M^3.
\]
The weak maximum principle applied with
$\phi(t)=\alpha M^2+CtM^3$ gives $u\leq CM^2$, hence
$|\nabla\Rm|\leq CM/\sqrt t$. This proves the displayed estimate for $k=1$;
the theorem's higher-derivative assertion is retained as stated.
\end{proof}

\begin{corollary}
\label{topping-maxpr-higher-derivatives}
\source{topping:page-0052}
\dcref{ch3:3.2}
\group{maximum-principle}
\lean{Topping.HasShiTowerOn,
Topping.mul_pow_le_affineBarrier_of_shiTower,
Topping.sqrt_le_div_of_shiTower}
Let $M>0$ and let $g(t)$ be a Ricci flow on a closed manifold $\M^n$ for
$t\in[0,\frac1M]$. For all $j,k\in\{0\}\cup\NN$, there is
$C=C(j,k,n)$ such that, if $|\Rm|\leq M$ on
$\M\times[0,\frac1M]$, then
\[
\left|\frac{\partial^j}{\partial t^j}\nabla^k\Rm\right|
\leq\frac{CM}{t^{j+\frac{k}{2}}}.
\]
\end{corollary}

\begin{proof}
\sketch Rescaling reduces the estimate to $t=1$, where $M\leq1$. The cases
$(j,k)=(1,0)$ and $(1,1)$ follow from the evolution equations for $\Rm$ and
$\nabla\Rm$ and Theorem~\ref{topping-maxpr-first-derivative-estimate}. For $j=1$ and $k>1$, induct on $k$ using
\[
\frac{\partial}{\partial t}\nabla^k\Rm
=\nabla\frac{\partial}{\partial t}\nabla^{k-1}\Rm
 +\nabla^{k-1}\Rm*\nabla\Rm.
\]
Each resulting $*$-product is controlled by the preceding derivative bounds;
iteration in $j$ gives the full estimate.
\end{proof}

\begin{theorem}
\label{topping-maxpr-shi-estimate}
\source{topping:page-0053}
\dcref{ch3:3.3}
\group{maximum-principle}
Let $g(t)$ be a Ricci flow on a boundaryless, not necessarily complete
manifold $U$ for $t\in[0,T]$. If $|\Rm|\leq M$ on $U\times[0,T]$, $p\in U$,
$r>0$, and $B_{g(0)}(p,r)\subset U$, then
\[
|\nabla\Rm(p,T)|^2\leq C(n)M^2
\left(\frac1{r^2}+\frac1T+M\right).
\]
\end{theorem}
